\documentclass[12pt]{article}
\usepackage[utf8]{inputenc}
\usepackage[spanish]{babel}
\usepackage{graphicx}
\usepackage{hyperref}
\usepackage{geometry}
\geometry{a4paper, margin=1in}

\title{DR - AGRAMONTE \\ \large Especificación Funcional del Proyecto}
\author{Mar Agramonte Campos}
\date{Mayo 2026}

\begin{document}

\maketitle

\section{Mapa de Navegación y Estructura}
A continuación se detalla la distribución de las secciones principales de la plataforma:

\begin{figure}[h]
\centering
\includegraphics[width=\textwidth]{Captura de pantalla 2026-05-13 222339.png}
\caption{Mapa de Secciones y Mockups de la Interfaz}
\end{figure}

\begin{itemize}
\item \textbf{Inicio:} Hero con imagen principal, navegación y footer.
\item \textbf{Sobre mí:} Perfil del médico, trayectoria e iconos de contacto.
\item \textbf{Servicios:} Consultas, chequeos y centros asociados.
\item \textbf{Atención:} Telemedicina, reserva de citas y recetas electrónicas.
\item \textbf{Reservas:} Calendario interactivo y selección de servicio.
\item \textbf{Testimonios:} Valoraciones 1-10 y reseñas de pacientes.
\end{itemize}

\section{Arquitectura Técnica y Datos}
Descripción del stack tecnológico y estructura de datos del sistema:

\begin{figure}[h]
\centering
\includegraphics[width=\textwidth]{Captura de pantalla 2026-05-13 222349.png}
\caption{Estructura de Base de Datos y Stack Tecnológico}
\end{figure}

\subsection{Estructura de la Ficha Médica}
La base de datos incluye campos para nombre, fecha de nacimiento, antecedentes, enfermedad actual, exámenes complementarios (analíticas, Rx, ecografía, TAC, RNM) y tratamiento.

\subsection{Stack Tecnológico}
El desarrollo utiliza HTML, CSS y JavaScript para el frontend; Java para el backend; y SQL para la gestión de base de datos. El testing se realiza mediante API REST.

\end{document}